% !PRINTED none
\begin{center}
\gesperrt{Sixteenth Chapter.}

\textbf{Numbers and Functionals of an Algebraic Field.}

\rule{4cm}{0.4pt}
\end{center}

\begin{center}
\S.\,132.

\gesperrt{Definition of the algebraic numbers.}
\end{center}

An algebraic equation
\[
F(x) = x^n + a_1 x^{n-1} + \cdots + a_{n-1} x + a_n = 0,
\]
whose coefficients $a_1, a_2, \ldots, a_n$ are rational numbers, we call for brevity a rational equation. It has, as we have proved in earlier chapters, always $n$ or fewer roots, but always at least one. How one can come as near as one pleases to each of these roots by rational numbers, say by decimal fractions or by continued fractions, if need be with the aid of the imaginary unit $i = \sqrt{-1}$, that is, how one can compute the values of the roots approximately, has been shown in the second Book of the first volume.

In the following considerations we shall not be concerned with these numerical values, but with the arithmetical laws to which these numbers are subject, and which can be derived from the definition itself and not from the numerical values. We therefore now place the following definition at the head:

\gesperrt{A number $\Theta$ which satisfies a rational equation}
\weq{1}{F(\Theta) = 0}
\gesperrt{is called an algebraic number.}

Every algebraic equation with rational coefficients furnishes us such algebraic numbers, which can therefore be given in arbitrary quantity. The question whether there are also numbers that are not algebraic will occupy us later.
